\section*{الفصل \ref{ch:b2:plant-life-cycles} --- دورات حياة النباتات اليابسة وتكاثرها}

\begin{solution}{exo:b2:plant-life-cycles:1}
\omterm{def:b2:plant-life-cycles:cycles}{الطور المشيجي}: الجيل \omterm{def:b2:meiosis-heredity:meiosis}{الفردي الصيغة} \omterm{def:b2:unicellular-diversity:unicellular}{العديد الخلايا}، النامي بالتقسم
الخيطي من بوغ، والذي يصنع الأعراس بالتقسم الخيطي. \omterm{def:b2:plant-life-cycles:cycles}{والطور البوغي}:
الجيل \omterm{def:b2:meiosis-heredity:meiosis}{الثنائي الصيغة}، النامي بالتقسم الخيطي من لاقحة، والذي يصنع
الأبواغ \omterm{def:b2:meiosis-heredity:meiosis}{بالتقسم المنصّف}. والبوغ: خلية فردية الصيغة مصنوعة \omterm{def:b2:meiosis-heredity:meiosis}{بالتقسم
المنصّف} تنمو دون اندماج. والعروس: خلية فردية الصيغة مصنوعة بالتقسم
الخيطي (عند النباتات) عليها أن تندمج بأخرى.
\end{solution}

\begin{solution}{exo:b2:plant-life-cycles:2}
خلايا الورقة $2n = 1260$: فالبوغ 630، وخلية \omterm{def:b2:plant-life-cycles:fern}{المشيجة} 630، والنطفة
630، والبويضة 630، واللاقحة 1260.
\end{solution}

\begin{solution}{exo:b2:plant-life-cycles:3}
الوسادة الخضراء: \omterm{def:b2:plant-life-cycles:cycles}{طور مشيجي}. والساق ذات المحفظة: \omterm{def:b2:plant-life-cycles:cycles}{طور بوغي}. وسعفة
\omterm{def:b2:plant-life-cycles:fern}{السرخس}: \omterm{def:b2:plant-life-cycles:cycles}{طور بوغي}. \omterm{def:b2:plant-life-cycles:fern}{والمشيجة}: \omterm{def:b2:plant-life-cycles:cycles}{طور مشيجي}. وحبة \omterm{def:b2:plant-life-cycles:heterospory}{اللقاح}: \omterm{def:b2:plant-life-cycles:cycles}{طور مشيجي} ذكري.
ومخزون غذاء البذرة (عند صنوبر): \omterm{def:b2:plant-life-cycles:cycles}{طور مشيجي} أنثوي.
وجنين البذرة: \omterm{def:b2:plant-life-cycles:cycles}{الطور البوغي} التالي.
\end{solution}

\begin{solution}{exo:b2:plant-life-cycles:4}
فردية الصيغة (\emph{Chlamydomonas}): \omterm{def:b2:meiosis-heredity:meiosis}{التقسم المنصّف} في اللاقحة،
حالًا. وثنائية الصيغة (الحيوانات، و\emph{Fucus}): \omterm{def:b2:meiosis-heredity:meiosis}{التقسم المنصّف}
يصنع الأعراس. وفردية--ثنائية الصيغة (الطحالب، والسراخس، والنباتات
البذرية، وكثير من الطحالب المائية): \omterm{def:b2:meiosis-heredity:meiosis}{التقسم المنصّف} يصنع الأبواغ في
الحوافظ البوغية \omterm{def:b2:plant-life-cycles:cycles}{للطور البوغي}.
\end{solution}

\begin{solution}{exo:b2:plant-life-cycles:5}
$0.03/10^{-4} = \qty{300}{s}$، أي خمس دقائق. وفي \qty{20}{min}،
تبلغ \qty{12}{cm}. فالإلقاح لا ينجح إلا بين نباتات تبعد بضعة
سنتيمترات، ولذلك تنمو الطحالب \omterm{def:b2:meiosis-heredity:vocabulary}{وسائد} كثيفة يكون فيها الجنسان (أو
العضوان) في متناول أحدهما الآخر.
\end{solution}

\begin{solution}{exo:b2:plant-life-cycles:6}
$v_s = 2\times(1.5\times 10^{-5})^2\times 1099\times 9.81/(9\times
1.8\times 10^{-5}) = \qty{0.030}{m/s}$، أي \qty{3}{cm/s}؛ والمسافة $uh/v_s
= 2\times 0.5/0.03 = \qty{33}{m}$. أما البذرة: $2\times 20/0.8 =
\qty{50}{m}$ --- فوحدة أثقل تُطلق من أعلى ولها جناح تقطع مسافة
مماثلة تقريبًا.
\end{solution}

\begin{solution}{exo:b2:plant-life-cycles:7}
$200\times 40\times 64 = \num{512000}$ بوغ في السعفة؛ ومنها $5.1$
\omterm{def:b2:plant-life-cycles:fern}{مشيجة}؛ أي $0.1$ \omterm{def:b2:plant-life-cycles:cycles}{طور بوغي} فتي في السعفة.
\end{solution}

\begin{solution}{exo:b2:plant-life-cycles:8}
البذرة \omterm{def:b2:plant-life-cycles:heterospory}{بييضة} مستبقاة: فعلى \omterm{def:b2:plant-life-cycles:cycles}{الطور المشيجي} الأنثوي أن يبقى على الأم،
محاطًا ومغذًّى، بينما على \omterm{def:b2:plant-life-cycles:cycles}{الطور المشيجي} الذكري أن يسافر إليه. وهذا
يستلزم صنفين من الأبواغ بمصيرين. أما البوغ الوحيد للنبات المتماثل
الأبواغ فلا بد من نثره لينمو طورًا مشيجيًا ثنائي الجنس حرًا؛ ولو
استُبقي لاستُبقيت معه الوظيفة الذكرية أيضًا ولفُرض الإلقاح الذاتي
على كل نبات، ولما بقي شيء يسافر.
\end{solution}

\begin{solution}{exo:b2:plant-life-cycles:9}
البوغ: $\tfrac43\pi(1.5\times 10^{-5})^3\times 1100 =
\qty{1.6e-11}{kg}$، أي \qty{16}{ng}؛ وطاقته $0.5\times 1.6\times
10^{-8}\times 2\times 10^{4} = \qty{1.6e-4}{J}$. والبذرة: \qty{5}{mg}،
وطاقتها $0.5\times 5\times 10^{-3}\times 2\times 10^{4} = \qty{50}{J}$ ---
أي أكثر بمقدار $3\times 10^{5}$ مرة. والبوغ لا يستطيع أن ينمي إلا حرشفة
صغيرة مركّبة ضوئيًا ولا شيء قبل الضوء؛ أما البذرة فتستطيع أن تنمي
جذرًا وساقًا في الظلام، وأن تنتظر موسمًا، وأن تنجو من أن تُؤكل أو
تجف.
\end{solution}

\begin{solution}{exo:b2:plant-life-cycles:10}
$10^{11}/(10^{6}\times 20) = \num{5000}$ حبة في المتر المكعب؛ والتدفق
$5000\times 10^{-7}\times 2 = \qty{1e-3}{}$ حبة في الثانية؛ أي نحو
\qty{1000}{s}، ربع ساعة، لكل حبة. فعلى المخروط أن يبقى مفتوحًا
مستقبِلًا ساعات إلى أيام، وذلك بالضبط حين تنثر المخاريط الذكرية ---
فالإلقاح اللقاحي موقوت إلى الأسبوع.
\end{solution}

\begin{solution}{exo:b2:plant-life-cycles:11}
تقنّع ثنائية الصيغة الطفرات المتنحية الضارة؛ والجسم \omterm{def:b2:meiosis-heredity:meiosis}{الثنائي الصيغة}
ذو النسيج الوعائي يستطيع أن يكون كبيرًا طويل العمر، والعلو يعين على
التقاط الضوء وعلى \omterm{thm:b2:plant-life-cycles:dispersal}{انتشار الأبواغ} معًا؛ والأبواغ المصنوعة عاليًا على
\omterm{def:b2:plant-life-cycles:cycles}{طور بوغي} تنتشر في الهواء، بينما على الأعراس أن تسبح. ويبقى الجيل
\omterm{def:b2:meiosis-heredity:meiosis}{الفردي الصيغة} لأن \omterm{def:b2:meiosis-heredity:meiosis}{التقسم المنصّف} والإلقاح يستلزمان طورًا فرديًا
الصيغة، ولأن حبة \omterm{def:b2:plant-life-cycles:heterospory}{اللقاح} والكيس الجنيني أصغر جسمين يحملان الأعراس
أحدهما إلى الآخر ويغذيان الجنين.
\end{solution}

\begin{solution}{exo:b2:plant-life-cycles:12}
بزرع الأبواغ وتتبعها تحت المجهر استطاع أن يثبت أن البوغ ينمو جسمًا
صغيرًا يحمل أعضاء جنسية، وأن الجنين ينشأ من البويضة الملقحة داخل
الرحم، وأن النبات الحامل للأبواغ ينمو منه، وأن الأعضاء نفسها موجودة
مختزلة في \omterm{def:b2:plant-life-cycles:heterospory}{بييضة} الصنوبر: أي التناوب بوصفه تعاقب أجسام وأعضاء. ولم
يستطع أن يعرف ما يميز الجيلين على مستوى الخلية. أما عدّ الصبغيات
فبيّن أن الجسمين يختلفان بعامل اثنين في عدد الصبغيات، وأن \omterm{def:b2:meiosis-heredity:meiosis}{التقسم
المنصّف} يقع عند تكوّن الأبواغ والإلقاح عند البويضة، فشرح بذلك لماذا
كان التناوب إلزاميًا.
\end{solution}

\begin{solution}{pb:b2:plant-life-cycles:1}
\textbf{1.} $10\times 200\times 40\times 64 = 5.1\times 10^{6}$
بوغ.
\textbf{2.} $\tfrac43\pi(1.5\times 10^{-5})^3\times 1100 =
\qty{1.6e-11}{kg}$ للبوغ الواحد؛ والمحصول $8\times 10^{-5}$ كغ، أي \qty{80}{mg}.
\textbf{3.} $v_s = 2\times 2.25\times 10^{-10}\times 1099\times
9.81/(1.62\times 10^{-4}) = \qty{0.030}{m/s}$.
\textbf{4.} $x = 1.5\times 0.5/0.030 = \qty{25}{m}$؛ وقرص مساحته
$\pi\times 25^2 = \qty{2000}{m^2}$؛ أي نحو \num{2600} بوغ في المتر
المربع.
\textbf{5.} التيارات الصاعدة ببضعة سنتيمترات في الثانية تفوق $v_s$،
فالبوغ الذي يعلق في اضطراب يبقى في الجو ساعات وينتقل مئات
الكيلومترات؛ وبوغ واحد يؤسس جمهرة إذا استطاعت مشيجته الإلقاح الذاتي
(فهي تحمل العضوين)، ولهذا كانت للجزر النائية فلورا سراخس غنية.
\textbf{6.} البوغ: $0.5\times 1.6\times 10^{-8}\times 2\times 10^{4} =
\qty{1.6e-4}{J}$؛ والمحصول: \qty{820}{J}؛ وبذرة صنوبر واحدة: \qty{50}{J} ---
أي إن محصول الأبواغ كله يعادل ست عشرة بذرة.
\textbf{7.} $5.1\times 10^{6}/2\times 10^{4} = 256$ \omterm{def:b2:plant-life-cycles:fern}{مشيجة}.
\textbf{8.} $1 - 0.8^{6} = 1 - 0.26 = 0.74$.
\textbf{9.} $256\times 0.74\times 0.5\times 0.02 = 1.9$ \omterm{def:b2:plant-life-cycles:fern}{سرخس} بالغ
في الموسم.
\textbf{10.} $30\times 1.9 \approx 57$: أي أكثر بكثير من البديل
الواحد اللازم لجمهرة مستقرة، فأرقام النجاة متفائلة --- ففي موئل
ممتلئ يموت معظم الأطوار البوغية الفتية من الازدحام والظل.
\textbf{11.} $10^{-3}/10^{-4} = \qty{10}{s}$. وسراخس كثيرة تنضج
مذاريها وأرحامها في أوقات مختلفة، أو تطلق هرمونًا (مولّد المذاري)
يجعل المشيجات الفتية المجاورة ذكرية، فتأتي النطاف من فرد آخر.
\textbf{12.} هي بسمك خلية واحدة، بلا جذور، بلا حماية، وتحتاج إلى
ماء سائل للإلقاح وتعيش بضعة أسابيع: فجل وفيات الدورة تقع عليها
(فبوغ من كل $2\times 10^{4}$ يصير \omterm{def:b2:plant-life-cycles:fern}{مشيجة}؛ وربع هذه لا يرى ليلة رطبة
أبدًا).
\textbf{13.} $v_s = 2\times 6.25\times 10^{-10}\times 599\times
9.81/(1.62\times 10^{-4}) = \qty{0.045}{m/s}$.
\textbf{14.} $3\times 15/0.045 = \qty{1000}{m}$.
\textbf{15.} $10^{11}/(2\times 10^{7}) = \num{5000}$ حبة في المتر المكعب.
\textbf{16.} $5000\times 10^{-7}\times 2 = \qty{1e-3}{s^{-1}}$: أي 3.6 في
الساعة؛ وتصل الأولى بعد نحو \qty{1000}{s}، أي ربع
ساعة.
\textbf{17.} على بعد خمسة كيلومترات مع الريح تكون السحابة قد انتشرت
جانبيًا وصاعدًا واستقر معظم الحبات (بمدى \qty{1}{km} لكل
\qty{15}{m} من الارتفاع)، فيكون التركيز أدنى برتب من القدر وقد
تنتظر \omterm{def:b2:plant-life-cycles:heterospory}{البييضة} أيامًا لحبة واحدة؛ \omterm{def:b2:plant-life-cycles:heterospory}{ولقاح} الشجرة نفسها يعطي في معظمه
بذورًا ملقحة ذاتيًا تجهض. فالإلقاح بالرياح يحتاج إلى سحابة كثيفة،
ومن ثم إلى غابات.
\textbf{18.} $10^{14}$ حبة من أجل $2\times 10^{6}$ \omterm{def:b2:plant-life-cycles:heterospory}{بييضة}: أي $5\times
10^{7}$ حبة لكل \omterm{def:b2:plant-life-cycles:heterospory}{بييضة}.
\textbf{19.} $2\times 10^{6}$ \omterm{def:b2:plant-life-cycles:heterospory}{بييضة} $\times 0.6\times 0.8 = 9.6\times
10^{5}$ بذرة؛ أي \qty{4.8}{kg}؛ و\qty{48}{MJ}.
\textbf{20.} $3\times 20/0.8 = \qty{75}{m}$: \omterm{def:b2:plant-life-cycles:heterospory}{فاللقاح} يقطع
كيلومترًا، والبذرة بضعة أطوال شجرة؛ فالمورّثات تسافر \omterm{def:b2:plant-life-cycles:heterospory}{باللقاح}،
والنبات بالبذرة.
\textbf{21.} $9.6\times 10^{5}\times 0.05\times 0.01 = 480$ خلفًا
بالغًا، في مقابل 57 عند \omterm{def:b2:plant-life-cycles:fern}{السرخس}.
\textbf{22.} الصنوبر: $4.8\times 10^{7}/480 = \qty{100}{kJ}$ لكل خلف
بالغ؛ \omterm{def:b2:plant-life-cycles:fern}{والسرخس}: $820/1.9 = \qty{430}{J}$ لكل خلف بالغ، أي أقل بمئتي
مرة.
\textbf{23.} كلاهما نصفه مادة جافة عند \qty{20}{kJ/g}، فعند معدل
قاعدي قدره مثلًا \qty{0.5}{mW} لكل غرام جاف يدوم كلاهما $10^{4}\,
\text{J/g}/(5\times 10^{-4}\,\text{W/g}) = 2\times 10^{7}$ ثانية، أي
نحو ثمانية أشهر: فمدة السكون ليست ما تشتريه طاقة البذرة. وإنما
تشتري ما يقع بعد الإنبات --- جذرًا وساقًا يُبنيان في الظلام من
المؤونة، قبل أسابيع من تغذي البادرة نفسها --- وتشتري، مع القشرة،
حمايةً في تلك الأثناء؛ أما \qty{1.6e-4}{J} التي للبوغ فتبني بضع
خلايا عليها أن تركّب ضوئيًا حالًا.
\textbf{24.} البذرة: إلقاح بلا ماء؛ وجنين مموَّن محمي ساكن يتأسس في
أماكن جافة أو موسمية؛ وانتشار بالأجنحة والثمار والحيوانات. أما
استراتيجية \omterm{def:b2:plant-life-cycles:fern}{السرخس} فتغلب في الموائل الرطبة المظللة المستقرة، وفي
استعمار أرض نائية أو حديثة الانكشاف، حيث تهم النواتج التكاثرية
الرخيصة الخفيفة الكثيرة أكثر من التموين.
\textbf{25.} \omterm{def:b2:plant-life-cycles:fern}{السرخس}: $5.1\times 10^{6}$ بوغ في الموسم مقابل 1.9 بالغ،
أي $2.7\times 10^{6}$ بوغ لكل خلف بالغ، و\qty{430}{J} لكل منها؛
والصنوبر: $9.6\times 10^{5}$ بذرة مقابل 480 بالغًا، أي \num{2000} بذرة
لكل خلف بالغ، و\qty{100}{kJ} لكل منها.
\end{solution}
